\documentclass[a4paper,12pt]{article}

\usepackage[T2A]{fontenc}
\usepackage[utf8]{inputenc}
\usepackage[russian]{babel}

\usepackage{amsmath, amssymb}
\usepackage{geometry}
\usepackage{fancyhdr}
\usepackage{enumitem}

\geometry{margin=2cm}
\setlength{\parskip}{6pt}
\setlength{\parindent}{0pt}

\pagestyle{fancy}
\fancyhf{}
\cfoot{\thepage}
\rhead{8 класс}

\begin{document}

\begin{center}
    {\Large \textbf{Графы. Начало.}}\\[0.5em]
\end{center}
    
\begin{enumerate}[label=\textbf{Задача \arabic*.}]
  \item Докажите, что не существует графа с пятью вершинами, степени которых равны 4, 4, 4, 4, 2.
  \item Докажите, что для всех n существует граф с 2n вершинами, степени которых равны $1, 1, 2, 2,\ldots, n,n$.
  \item
  Сколько существует графов с 10 вершинами, степень каждой из которых равна 9?
  \item В стране любые два города связаны авиалинией одной из двух авиакомпаний. Докажите, что можно закрыть одну из авиакомпаний так, что по–прежнему от любого города можно будет добраться до любого другого.
  \item
  \begin{enumerate}[label=\textbf{\alph*)}]
    \item Докажите, что в графе число вершин нечётной степени чётно.
    \item Какое максимальное число ребер может быть в графе с n вершинами?
    \item Какое максимальное число ребер может быть в графе с n вершинами и k компонентами связности (считайте, что размер i-ой компоненты - $p_i$ и $\displaystyle \sum_{i = 1}^{k}p_i = n$ (то же самое, что и $p_1 + p_2 + p_3 + \ldots + p_k = n$))?
  \end{enumerate}

  \item Докажите, что для любого графа \(G\) можно приписать каждой вершине \(u\) число \(d(u)\) так, что для любых вершин \(u,v\) числа \(d(u)\) и \(d(v)\) имеют общий делитель, отличный от \(1\), тогда и только тогда, когда в графе \(G\) есть ребро \(uv\).
  
  \item Докажите, что если в графе с \(n\) вершинами минимальная степень \(\delta(G) > \frac{n-1}{2}\), то граф \(G\) связен ($\delta(G)$ - минимальная степень вершины в графе G).

\item В связном графе нашлись два простых пути максимальной длины. Докажите, что у них есть хотя бы одна общая вершина.

\item Пусть T — дерево с хотя бы 3 вершинами. Докажите, что T имеет хотя бы одну из двух следующих конфигураций:
(*) два листа, смежных с одной и той же вершиной;
(**) лист, смежный с вершиной степени 2.

\item В выпуклом n-угольнике провели максимально возможное количество диагоналей, не имеющих общих внутренних точек, которые разбили его на треугольники (это называется триангулировать многоугольник).
а) Докажите, что получилось ровно n - 2 треугольника.
б) Пусть каждый треугольник — вершина графа, а две вершины смежны, если треугольники имеют общую сторону. Докажите, что получилось дерево.

\item В вершинах графа расставлены числа ±1 таким образом, что для каждой вершины произведение чисел в вершинах, соседних с ней, отрицательно. Чему может быть равно произведение чисел во всех вершинах?

\item Вершины дерева T покрашены в красный и синий цвет так, что любые две смежные вершины разноцветны, и красных не меньше, чем синих. Докажите, что у T есть красный лист.
\item Может ли сумма попарных расстояний между вершинами 25-вершинного дерева быть равна 1225?

\item Среди n рыцарей каждые двое – либо друзья, либо враги. У каждого из рыцарей ровно три врага, причём враги его друзей являются его врагами. При каких n такое возможно?

\item На авиаслёте собрались 49 пилотов.
Каждый пилот общается как минимум с 25 другими участниками — они знакомы и часто летают вместе. Докажите, что всех пилотов можно разбить на небольшие экипажи — по 2 или 3 человека — так, чтобы в каждом экипаже все его члены были знакомы друг с другом.

\item Среди n бизнесменов каждые двое – либо конкрутенты, либо партнёры. У каждого из бизнесмена ровно три конкурента, причём конкуренты его партнёра являются его конкурентами. При каких n такое возможно?

    
\end{enumerate}
\end{document}
